
\section{Conclusion}
본 논문에서는 ``LLM은 네트워크 설정 파일을 얼마나 정확하게 해석하고 분석할 수 있는가?''라는 질문에 답하기 위해 벤치마크 NetConfigQA~2.0, 평가 지표 TA-Acc, 도구 통합 Multi-Agent 시스템 NetAlly를 제안하였다. 이중 경로 파이프라인은 설정 파일만으로 QA를 자동 생성하였으며, 4개 토폴로지(10-40 노드, 9,473 QA)에 코드 수정 없이 적용되었고 3중 GT 검증에서 98-99.9\%의 일치율을 보였다. 6개 LLM을 평가한 결과, 사실 추출(L1-L3)에서는 모델 간 뚜렷한 성능 차이가 나타났으나 경로 추론과 장애 분석(L4-L5)에서는 모든 모델이 시뮬레이션 장벽에 직면하였다. 3-way 비교에서 MAS 구조만으로는 L4/L5가 개선되지 않았으나 네트워크 도구 통합 시 유의미한 변화가 관찰되어, 성능 향상의 주된 원천이 도구 통합에 있음을 확인하였다. TA-Acc는 BERTScore와 Exact Match가 구분하지 못하는 모델 간 차이를 식별하였다.

우리는 향후 연구로 L6 트러블슈팅 확장, 멀티 벤더(Juniper, Arista) 지원, 설정 생성 태스크 평가, 외부 벤치마크와의 교차 적용을 계획한다.
